% ======================================================================
% content/folien.tex  (OHNE Praeambel)
% Stand: 2026-01-21 | Version: 2.8.0
% Aktualisiert: Pfade angepasst an aktuelle Repo-Struktur (apps/*)
% Farbschema: ArduinoTeal + RaspberryRed
%
% Annahmen zur Ordnerstruktur:
% - Datei liegt in:   ./latex/content/folien.tex
% - Bilder liegen in: ./latex/images/
%
% Konsequenz:
% - KEIN "images/..." in \includegraphics, sondern nur Dateiname.
% - \graphicspath ist im praesentation.tex bereits auf ../images/ gesetzt.
% ======================================================================

% --- Folie 1: Titel (Cover als Hintergrund) ----------------------------
% In praesentation.tex definiert.

% --- Folie 2: Hardware-Prototyp ----------------------------------------
\begin{frame}{Hardware-Prototyp v2}
\framesubtitle{10 Taster \textbullet\ 10 LEDs \textbullet\ Aufger\"{a}umtes Layout}

\begin{columns}[c,onlytextwidth]
\column{0.48\textwidth}
\centering
\fbox{\includegraphics[width=0.95\linewidth,height=4.5cm,keepaspectratio]{prototyp-v2-front.png}}
\vspace{0.1cm}

{\small\textbf{Best\"{u}ckungsseite} -- ESP32 + ICs}

\column{0.48\textwidth}
\centering
\fbox{\includegraphics[width=0.95\linewidth,height=4cm,keepaspectratio]{schaltplan.pdf}}
\vspace{0.1cm}

{\small\textbf{KiCad} -- Schaltplan}
\end{columns}

\vspace{0.2cm}
\highlight{Komponenten:} ESP32-S3 XIAO \textbullet\ 2$\times$ 74HC595 \textbullet\ 2$\times$ CD4021BE \textbullet\ 10 LEDs \textbullet\ 10 Taster
\end{frame}

% --- Folie 2.5: Blockdiagramm Prototyp ---------------------------------
\begin{frame}{Schieberegister-Architektur}
\framesubtitle{Prototyp 10$\times$ -- 6 GPIO-Pins f\"{u}r 20 IOs}

\begin{center}
\includegraphics[width=0.85\linewidth,height=0.5\textheight,keepaspectratio]{prototyp-10x.png}%

\end{center}

\vspace{-0.2em}

\begin{columns}[T]
\column{0.48\textwidth}
\small
\textbf{\textcolor{ArduinoTeal}{Komponenten}}
\begin{itemize}\setlength{\itemsep}{0pt}
  \item \textcolor{ArduinoTeal}{ESP32-S3 XIAO} -- Mikrocontroller
  \item \textcolor{RaspberryRed}{2$\times$ 74HC595} -- LED-Ausgabe (Active-High)
  \item \textcolor{SuccessGreen}{2$\times$ CD4021BE} -- Taster-Eingabe (Active-Low)
\end{itemize}

\column{0.48\textwidth}
\small
\textbf{\textcolor{ArduinoTeal}{Signale (SPI-\"{a}hnlich)}}
\begin{itemize}\setlength{\itemsep}{0pt}
  \item \textcolor{SignalData}{Data}: MOSI $\rightarrow$ 74HC595 SER, MISO $\leftarrow$ CD4021 Q8
  \item \textcolor{SignalClock}{Clock}: SCK gemeinsam (SRCLK/CLK)
  \item \textcolor{SignalLatch}{Strobes}: RCK (595 Latch), P/S (4021 Load/Shift)
\end{itemize}
\end{columns}
\end{frame}

% --- Folie 2.7: Prototyp-Evolution -------------------------------------
\begin{frame}{Prototyp-Evolution}
\framesubtitle{v1 $\rightarrow$ v2: Von der Idee zum aufger\"{a}umten Layout}

\begin{columns}[c,onlytextwidth]
\column{0.48\textwidth}
\centering
\fbox{\includegraphics[width=0.95\linewidth,height=2.5cm,keepaspectratio]{prototyp-v1-front.png}}%

\vspace{0.1cm}

{\small\textbf{v1 Best\"{u}ckung} -- Sockel ohne ICs}

\vspace{0.2cm}
\fbox{\includegraphics[width=0.95\linewidth,height=2.5cm,keepaspectratio]{prototyp-v1-back.png}}%

\vspace{0.1cm}

{\small\textbf{v1 L\"{o}tseite}}

\column{0.48\textwidth}
\centering
\fbox{\includegraphics[width=0.95\linewidth,height=2.5cm,keepaspectratio]{prototyp-v2-front.png}}
\vspace{0.1cm}

{\small\textbf{v2 Best\"{u}ckung}}

\vspace{0.2cm}

\fbox{\includegraphics[width=0.95\linewidth,height=2.5cm,keepaspectratio]{prototyp-v2-back.png}}%

\vspace{0.1cm}

{\small\textbf{v2 L\"{o}tseite}}
\end{columns}
\end{frame}

% --- Folie 3: Web-Dashboard --------------------------------------------
\begin{frame}{Web-Dashboard}
\framesubtitle{React-Frontend \textbullet\ WebSocket \textbullet\ Multi-Device}

\begin{center}
\fbox{\includegraphics[width=0.85\linewidth,height=5cm,keepaspectratio]{web-dashboard2.png}}%

\end{center}

\vspace{0.2cm}
\begin{columns}[T,onlytextwidth]
\column{0.24\textwidth}\centering
\textbf{Zugriff}\\{\small \texttt{rover.local:8081}}

\column{0.24\textwidth}\centering
\textbf{Multi-Device}\\{\small iMac, iPad, iPhone}

\column{0.24\textwidth}\centering
\textbf{React + Vite}\\{\small TypeScript, HMR}

\column{0.24\textwidth}\centering
\textbf{Synchron}\\{\small WebSocket Broadcast}
\end{columns}
\end{frame}

% --- Folie 4: Was ist das? ---------------------------------------------
\begin{frame}{Was ist das Selection Panel?}
\framesubtitle{Physische Auswahl $\rightarrow$ Digitale Reaktion}

\begin{columns}[T,onlytextwidth]
\column{0.48\textwidth}
\textbf{Anwendungsfall}
\begin{itemize}
  \item Prototyp: 10 physische Taster
  \item Vollausbau: 100 Taster geplant
  \item LED zeigt aktive Auswahl an
  \item Web-Dashboard visualisiert Status
\end{itemize}

\column{0.48\textwidth}
\textbf{Policy: \glqq{}LowestIdWins\grqq{}}
\begin{itemize}
  \item Neuer Tastendruck $\rightarrow$ \textbf{sofortiger Wechsel}
  \item Bei Mehrfachdruck: niedrigste ID gewinnt
  \item \textbf{One-hot LED:} Immer max.\ 1 LED aktiv
  \item Selection $=0$ $\rightarrow$ alle LEDs aus
\end{itemize}
\end{columns}

\vspace{0.3cm}
\highlight{Merksatz:} Bei gleichzeitigem Druck gewinnt die niedrigste ID -- deterministisch und vorhersagbar.
\end{frame}

% --- Folie 5: Systemarchitektur ----------------------------------------
\begin{frame}{Systemarchitektur -- \"{U}berblick}
\framesubtitle{Drei Ebenen: Controller \textbullet\ Bridge \textbullet\ Dashboard}

\begin{center}
\begin{tikzpicture}[
  box/.style={draw, rounded corners, minimum width=3.2cm, minimum height=1.2cm, align=center},
  arrow/.style={->, thick, >=stealth}
]
  \node[box, fill=ArduinoTeal!20] (esp) at (0,0) {\textbf{ESP32-S3}\\Firmware};
  \node[box, fill=RaspberryRed!20] (pi) at (5,0) {\textbf{Raspberry Pi 5}\\Bridge};
  \node[box, fill=SuccessGreen!20] (browser) at (10,0) {\textbf{Web-Browser}\\React};

  \draw[arrow, <->] (esp) -- node[above, font=\footnotesize] {USB-CDC} node[below, font=\footnotesize] {COBS + CRC16} (pi);
  \draw[arrow, <->] (pi) -- node[above, font=\footnotesize] {WebSocket} node[below, font=\footnotesize] {:8081 / JSON} (browser);
\end{tikzpicture}
\end{center}

\vspace{0.3cm}
\begin{table}
\centering
\small
\begin{tabular}{@{}ll@{}}
\toprule
\textbf{Komponente} & \textbf{Aufgabe} \\
\midrule
\textcolor{ArduinoTeal}{ESP32-S3 XIAO} & Taster scannen, LEDs steuern (10$\times$ Prototyp / 100$\times$ Vollausbau) \\
\textcolor{RaspberryRed}{Raspberry Pi 5} & Protokoll-Bridge (\texttt{apps/bridge/bridge.py}), WebSocket-Server \\
\textcolor{SuccessGreen}{Browser (React)} & Auswahl visualisieren, WebSocket-Client \\
\bottomrule
\end{tabular}
\end{table}
\end{frame}

% --- Folie 6: Hardware + CD4021 Timing ---------------------------------
\begin{frame}{Hardware -- Schieberegister}
\framesubtitle{Skalierbar: 6 GPIO-Pins f\"{u}r $n \times 8$ IOs}

\begin{columns}[T,onlytextwidth]
\column{0.50\textwidth}
\textbf{Warum Schieberegister?}
\begin{itemize}
  \item Prototyp: 20 IOs mit \textbf{6 GPIO-Pins}
  \item Vollausbau: $13 \times 8 = 104$ Kan\"{a}le (100 genutzt)
  \item Kaskadierbar, g\"{u}nstig und robust
\end{itemize}

\vspace{0.2cm}
\textbf{Pinbelegung ESP32-S3 (gemeinsamer SPI)}
\begin{table}
\footnotesize
\begin{tabular}{@{}lll@{}}
\toprule
\textbf{Signal} & \textbf{Pin} & \textbf{Funktion} \\
\midrule
SCK  & D8  & SPI Clock (gemeinsam) \\
MOSI & D10 & 74HC595 SER \\
MISO & D9  & CD4021BE Q8 \\
RCK  & D0  & 74HC595 Latch \\
OE   & D2  & 74HC595 Enable (Active-Low) \\
P/S  & D1  & CD4021BE Parallel/Shift \\
\bottomrule
\end{tabular}
\end{table}

\column{0.50\textwidth}
\centering
\includegraphics[width=0.98\linewidth,height=4.5cm,keepaspectratio]{CD4021-timing.png}%


\vspace{0.1cm}
{\small\textbf{CD4021BE Timing-Diagramm}}
\end{columns}
\end{frame}

% --- Folie 7: Firmware -------------------------------------------------
\begin{frame}{Firmware -- Modulare FreeRTOS-Architektur}
\framesubtitle{Schichtentrennung: HAL \textbullet\ Drivers \textbullet\ Logic \textbullet\ App}

\begin{columns}[T,onlytextwidth]
\column{0.50\textwidth}
\textbf{\textcolor{ArduinoTeal}{IO-Task} (Core 0, Prio 1)}
\begin{itemize}
  \item Polling alle $10\,\mathrm{ms}$ ($100\,\mathrm{Hz}$)
  \item Debouncing pro Taste ($50\,\mathrm{ms}$)
  \item Selection via IoPipeline (LowestIdWins)
\end{itemize}

\vspace{0.3cm}
\textbf{\textcolor{ArduinoDark}{Serial-Task} (Core 1, Prio 1)}
\begin{itemize}
  \item COBS-Framing (0x00 Delimiter)
  \item CRC-16 CCITT-FALSE
  \item ACK f\"{u}r Host-Kommandos (optional via Flag)
\end{itemize}

\column{0.50\textwidth}
\textbf{Schichtenmodell}
\begin{table}
\small
\begin{tabular}{@{}ll@{}}
\toprule
\textbf{Schicht} & \textbf{Inhalt} \\
\midrule
\textcolor{ArduinoTeal}{app/} & io\_task, serial\_task \\
\textcolor{ArduinoDark}{logic/} & debounce, selection, io\_pipeline \\
\textcolor{RaspberryRed}{drivers/} & cd4021, hc595 \\
\textcolor{SuccessGreen}{hal/} & spi\_bus + Mutex \\
\bottomrule
\end{tabular}
\end{table}

\vspace{0.2cm}
\textbf{SPI-Konfiguration}
\begin{itemize}
  \item CD4021BE: MODE1, $500\,\mathrm{kHz}$
  \item 74HC595: MODE0, $1\,\mathrm{MHz}$
  \item Mutex f\"{u}r Thread-Safety
\end{itemize}
\end{columns}
\end{frame}

% --- Folie 8: Protokoll ------------------------------------------------
\begin{frame}{Serial-Protokoll}
\framesubtitle{ESP32 $\leftrightarrow$ Raspberry Pi \"{u}ber USB-CDC}

\begin{columns}[T,onlytextwidth]
\column{0.48\textwidth}
\textbf{Frame-Aufbau (COBS-kodiert)}
\begin{table}
\small
\begin{tabular}{@{}ll@{}}
\toprule
\textbf{Feld} & \textbf{Gr\"{o}\ss e} \\
\midrule
Magic ('SP') & 2 Bytes \\
Version & 2 Bytes \\
Type & 1 Byte \\
Flags & 1 Byte \\
Sequence & 2 Bytes (LE) \\
Length & 2 Bytes (LE) \\
Payload & 0--256 Bytes \\
CRC-16 & 2 Bytes (LE) \\
\bottomrule
\end{tabular}
\end{table}

\vspace{0.2cm}
\textbf{Header:} 10 Bytes + Payload + 2 Bytes CRC

\column{0.48\textwidth}
\textbf{Nachrichtentypen (Firmware)}
\begin{itemize}\setlength{\itemsep}{1pt}
  \item \texttt{0x01 HELLO} -- Konfiguration (10 Bytes)
  \item \texttt{0x12 SELECTION\_STATE} -- Kanalwechsel (8 Bytes)
  \item \texttt{0x20 CMD\_SET\_SELECTION} -- Host-Befehl (4 Bytes)
  \item \texttt{0x30 ACK} -- Best\"{a}tigung (0 Bytes)
\end{itemize}

\vspace{0.2cm}
\textbf{ACK/Retry (Host $\rightarrow$ ESP32)}
\begin{itemize}
  \item Jede TX bekommt eine \texttt{seq} (1..65535, 0 wird \"{u}bersprungen)
  \item Parameter: \texttt{--ack-timeout} (typ.\ $0.6\,\mathrm{s}$), \texttt{--retries} (typ.\ 4)
  \item Bei Timeout: \texttt{warn} (code=20), bei Fail: \texttt{error} (code=21)
\end{itemize}
\end{columns}
\end{frame}

% --- Folie 9: Server ---------------------------------------------------
\begin{frame}[fragile]{Raspberry Pi Bridge}
\framesubtitle{Python 3 + aiohttp + pyserial-asyncio}

\textbf{Technologie-Stack:}
\begin{itemize}
  \item \textbf{aiohttp:} Async HTTP/WebSocket-Server
  \item \textbf{pyserial-asyncio:} Async Serial I/O
  \item \textbf{COBS/CRC:} Frame-Kodierung/-Pr\"{u}fung
  \item \textbf{systemd:} Autostart als Service (optional mit \texttt{flock})
\end{itemize}

\vspace{0.3cm}
\textbf{Datenfluss bei Tastendruck:}
\begin{enumerate}
  \item \textcolor{ArduinoTeal}{ESP32} sendet \texttt{SELECTION\_STATE} (z.\,B.\ \texttt{selected=3}, \texttt{source=local\_btn})
  \item \textcolor{RaspberryRed}{Bridge} speichert State, broadcastet JSON an alle Browser
  \item \textcolor{SuccessGreen}{Browser} empf\"{a}ngt \verb|{"type":"selection","selected":3,"source":"local_btn"}|
\end{enumerate}

\vspace{0.3cm}
\textbf{Datenfluss bei Browser-Klick:}
\begin{enumerate}
  \item \textcolor{SuccessGreen}{Browser} sendet \verb|{"type":"setSelection","selected":5}|
  \item \textcolor{RaspberryRed}{Bridge} sendet \texttt{CMD\_SET\_SELECTION} mit \texttt{seq} an ESP32
  \item \textcolor{ArduinoTeal}{ESP32} antwortet (optional) mit \texttt{ACK}, danach \texttt{SELECTION\_STATE}
  \item \textcolor{RaspberryRed}{Bridge} broadcastet neuen State an alle Browser
\end{enumerate}
\end{frame}

% --- Folie 10: Dashboard -----------------------------------------------
\begin{frame}{Dashboard -- Implementierung}
\framesubtitle{React + TypeScript \textbullet\ WebSocket \textbullet\ Statusleiste}

\begin{columns}[T,onlytextwidth]
\column{0.50\textwidth}
\textbf{Features}
\begin{itemize}
  \item \textbf{React + Vite:} HMR, TypeScript
  \item \textbf{WebSocket:} Auto-Reconnect (z.\,B.\ $1\,\mathrm{s}$)
  \item \textbf{Build (PROD):} Output nach \texttt{apps/dashboard/dist/}
  \item \textbf{Dev-Proxy:} zu \texttt{rover.local:8081}
\end{itemize}

\vspace{0.3cm}
\textbf{Farbschema (10 Kan\"{a}le)}
\begin{itemize}
  \item Wei\ss\ (1--2), Blau (3--4)
  \item Rot (5--6), Gelb (7--8)
  \item Gr\"{u}n (9--10), zyklisch erweiterbar
\end{itemize}

\column{0.50\textwidth}
\textbf{Statusleiste-Anzeige}
\begin{table}
\small
\begin{tabular}{@{}ll@{}}
\toprule
\textbf{Feld} & \textbf{Bedeutung} \\
\midrule
IN  & Eingangsbits (Active-Low) \\
OUT & Ausgangsbits (Active-High) \\
sel & Aktuell gew\"{a}hlter Kanal \\
src & Quelle (local\_btn / host\_cmd) \\
tMs & Timestamp (ms seit Boot) \\
\bottomrule
\end{tabular}
\end{table}

\vspace{0.2cm}
\textbf{Bedienung}
\begin{itemize}
  \item Klick auf Kanal-Tile $\rightarrow$ \texttt{setSelection}
  \item \glqq{}Alle aus\grqq{} $\rightarrow$ \texttt{selected=0}
\end{itemize}
\end{columns}
\end{frame}

% --- Folie 11: Tests ---------------------------------------------------
\begin{frame}{Getestete Funktionen}
\framesubtitle{End-to-End verifiziert -- Prototyp voll funktionsf\"{a}hig}

\begin{columns}[T,onlytextwidth]
\column{0.48\textwidth}
\textbf{\textcolor{ArduinoTeal}{Firmware (ESP32)}}
\begin{itemize}
  \item Modulare Architektur (HAL/Drivers/Logic/App)
  \item SPI-Bus mit Mutex (Thread-safe)
  \item Debouncing: $50\,\mathrm{ms}$ stabil
  \item Selection: One-hot LED-Ansteuerung (LowestIdWins)
  \item COBS/CRC16 Frame-Handling
\end{itemize}

\vspace{0.3cm}
\textbf{\textcolor{RaspberryRed}{Bridge (Raspberry Pi)}}
\begin{itemize}
  \item Async Serial (pyserial-asyncio)
  \item ACK/Retry f\"{u}r Host-Kommandos (parametrierbar)
  \item WebSocket Broadcast funktioniert
  \item State-Recovery bei Reconnect
\end{itemize}

\column{0.48\textwidth}
\textbf{\textcolor{SuccessGreen}{Dashboard (React)}}
\begin{itemize}
  \item 10-Kanal-Visualisierung
  \item WebSocket-Verbindung stabil
  \item Statusleiste zeigt Live-Daten
  \item Build-Pipeline (Vite $\rightarrow$ \texttt{apps/dashboard/dist/})
\end{itemize}

\vspace{0.3cm}
\textbf{\textcolor{SuccessGreen}{Hardware-Prototyp}}
\begin{itemize}
  \item 10 Taster funktionsf\"{a}hig
  \item 10 LEDs funktionsf\"{a}hig
  \item Bit-Mapping korrekt (MSB-first)
  \item CD4021BE + 74HC595 stabil
\end{itemize}
\end{columns}
\end{frame}

% --- Folie 12: Demo (ohne Embedded-Multimedia) --------------------------
\begin{frame}{Demo (End-to-End)}
\framesubtitle{Tasterdruck $\rightarrow$ LED \& WebSocket-Update}

\begin{center}
\fbox{\includegraphics[width=11cm,height=4cm,keepaspectratio]{video-thumbnail.png}}%

\end{center}

\vspace{0.2cm}
\begin{enumerate}\small
  \item Button $k$ wird gedr\"{u}ckt $\rightarrow$ \textcolor{ArduinoTeal}{ESP32} setzt \texttt{selected=k} (LowestIdWins).
  \item \textcolor{ArduinoTeal}{ESP32} sendet \texttt{SELECTION\_STATE} \"{u}ber USB-CDC (COBS + CRC16).
  \item \textcolor{RaspberryRed}{Bridge} broadcastet JSON \"{u}ber \texttt{/ws} an alle Browser.
\end{enumerate}
\end{frame}

% --- Folie 13: Live-Demo -----------------------------------------------
\begin{frame}[fragile]{Live-Demo}
\framesubtitle{Hardware-Prototyp + WebSocket-Test}

\textbf{Dashboard \"{o}ffnen (alle Ger\"{a}te):}
\begin{verbatim}
http://rover.local:8081/
\end{verbatim}

\vspace{0.2cm}
\textbf{WebSocket-Test mit wscat:}
\begin{verbatim}
wscat -c ws://rover.local:8081/ws
> {"type":"setSelection","selected":5}
\end{verbatim}

\vspace{0.2cm}
\textbf{Service-Status pr\"{u}fen:}
\begin{verbatim}
sudo systemctl status embedded-bridge.service
\end{verbatim}

\vspace{0.2cm}
\textbf{Serial-Port pr\"{u}fen:}
\begin{verbatim}
ls -la /dev/serial/by-id/ | grep Espressif
\end{verbatim}

\vspace{0.2cm}
\highlight{Hardware:} Taster dr\"{u}cken $\rightarrow$ LED leuchtet, Status auf allen Browsern synchron.
\end{frame}

% --- Folie 14: Ausblick + Vollausbau -----------------------------------
\begin{frame}{Ausblick: Vollausbau 100$\times$}
\framesubtitle{Skalierung von 10 auf 100 Kan\"{a}le}

\begin{columns}[c,onlytextwidth]
\column{0.55\textwidth}
\centering
\includegraphics[width=0.98\linewidth,height=5.5cm,keepaspectratio]{vollausbau-100x.png}%


\column{0.45\textwidth}
\textbf{\"{A}nderungen f\"{u}r 100$\times$}
\begin{itemize}
  \item 13$\times$ 74HC595 (statt 2$\times$)
  \item 13$\times$ CD4021BE (statt 2$\times$)
  \item \texttt{BTN\_COUNT = 100}
  \item \texttt{LED\_COUNT = 100}
  \item Gleiche Firmware/Bridge (nur Konfiguration)
\end{itemize}

\vspace{0.3cm}
\textbf{St\"{u}ckliste (100$\times$)}
\begin{table}
\footnotesize
\begin{tabular}{@{}lr@{}}
\toprule
\textbf{Bauteil} & \textbf{Anzahl} \\
\midrule
74HC595 & 13$\times$ \\
CD4021BE & 13$\times$ \\
LED $5\,\mathrm{mm}$ & 100$\times$ \\
Taster & 100$\times$ \\
R $330\,\Omega$ -- $3\,\mathrm{k\Omega}$ (LED) & 100$\times$ \\
R $10\,\mathrm{k\Omega}$ (Pull-up) & 100$\times$ \\
C $100\,\mathrm{nF}$ & 26$\times$ \\
\bottomrule
\end{tabular}
\end{table}
\end{columns}
\end{frame}

% --- Folie 15: Zusammenfassung -----------------------------------------
\begin{frame}{Zusammenfassung}
\framesubtitle{Modulare Architektur, skalierbar und robust}

\begin{table}
\centering
\small
\begin{tabular}{@{}lll@{}}
\toprule
\textbf{Phase} & \textbf{Inhalt} & \textbf{Status} \\
\midrule
1 & Infrastruktur \& Toolchain & \textcolor{SuccessGreen}{$\checkmark$} \\
2 & Hardware-Prototyp (10$\times$) & \textcolor{SuccessGreen}{$\checkmark$} \\
3 & ESP32 Firmware (FreeRTOS) & \textcolor{SuccessGreen}{$\checkmark$} \\
4 & Serial-Protokoll (COBS + CRC16) & \textcolor{SuccessGreen}{$\checkmark$} \\
5 & Raspberry Pi Bridge & \textcolor{SuccessGreen}{$\checkmark$} \\
6 & React-Frontend & \textcolor{SuccessGreen}{$\checkmark$} \\
7 & Hardware-Integration (10$\times$) & \textcolor{SuccessGreen}{$\checkmark$} \\
8 & Hardware-Skalierung (100$\times$) & \textcolor{RaspberryRed}{$\blacktriangleleft$} N\"{a}chste Phase \\
\bottomrule
\end{tabular}
\end{table}

\vspace{0.3cm}
\highlight{Kernbotschaft:} ESP32 f\"{u}r Echtzeit-I/O, Pi als Protokoll-Bridge, Browser f\"{u}r UI. Skalierung 10 $\rightarrow$ 100 durch Konfiguration.

\vspace{0.2cm}
\textbf{\"{U}bertragbarkeit:} Voting-Systeme, Quiz-Panels, Museumsinstallationen.
\end{frame}

% --- Folie 16: Credits -------------------------------------------------
\begin{frame}{Credits}
\framesubtitle{Projektbeteiligte und Werkzeuge}

\begin{center}
\begin{tikzpicture}[
  box/.style={draw=ArduinoTeal, rounded corners=8pt, minimum width=10cm, minimum height=1.2cm, align=center, fill=ArduinoTeal!5},
  label/.style={font=\small\bfseries, text=ArduinoDark}
]
  \node[box, draw=RaspberryRed, fill=RaspberryRed!5] (ki) at (0, 0.8) {%
    \textcolor{RaspberryRed}{\textbf{Claude Opus 4.5}} \textbullet\ \textcolor{RaspberryRed}{\textbf{ChatGPT 5.2}}\\[0.1cm]
    {\small\color{MutedText}Text- und Code-Vorschl\"{a}ge}
  };
  \node[label, anchor=east] at (-5.5, 0.8) {KI-Unterst\"{u}tzung};

  \node[box, draw=SuccessGreen, fill=SuccessGreen!5] (hinweis) at (0, -1.0) {%
    {\small\color{DarkText}Technische Pr\"{u}fung durch den Autor}
  };
  \node[label, anchor=east] at (-5.5, -1.0) {Qualit\"{a}tssicherung};
\end{tikzpicture}
\end{center}
\end{frame}
